% !PRINTED 657
\weberpage{657}

In diesem Falle durchl\"auft nun $\xi$ nach \S.\,172, III ein volles
System durch $q$ nicht theilbarer Reste nach dem Modul $m_1$, und
wir k\"onnen in (10) unter $\xi'$ die kleinste positive L\"osung der
Congruenz
\[
\xi\xi' \equiv 1\ (\mathrm{mod}\ m_1)
\]
verstehen, weil. wenn wir die Exponenten $\xi'$ in (10) nach dem % [source sic: full stop after ``weil'' as printed; a comma is grammatically expected]
Modul $m_1$ ver\"andern, nur $m^{\text{te}}$ Potenzen der
Primfactoren $\mathfrak{p}$ hinzutreten.

Dann k\"onnen wir auf das Product (10) das \gesperrt{Kummer'sche}
Theorem [\S.\,174, (16)] anwenden, indem wir dort $m$ durch $m_1$
ersetzen, und erhalten, wenn wir unter $r_1 = r^{m_2}$ eine
$m_1^{\text{te}}$ Einheitswurzel, unter den $\eta$ gewisse Perioden
der $p^{\text{ten}}$ Einheitswurzeln verstehen:
\[
\prod^{\xi} \mathfrak{p}_\xi^{\xi'} = (r_1, \eta)^{m_1},
\]
und folglich
\weq{11}{(\prod^{\xi} \mathfrak{p}_\xi^{\xi'})^{hm_2} = (r_1, \eta)^{hm}.}
Demnach haben wir das Theorem:

6. \gesperrt{Ist $p$ eine Primzahl und $p-1$ durch $q$, oder,
wenn $q=2$ ist, durch 4 theilbar, so l\"asst sich der positive
Exponent $h$ so bestimmen, dass die in $\Omega_m$ enthaltene Zahl}
\[
\omega (r_1, \eta)^{-hm}
\]
\gesperrt{die Primfactoren von $p$ nur zu solchen Potenzen
enth\"alt, deren Exponent durch $m$ theilbar ist.}

Es ist jetzt auch nicht mehr n\"othig, den Satz 5. von dem
Satze 6. zu unterscheiden; denn f\"ur $m_1 = 2$ geht 6. in 5.
\"uber, weil dann $r_1 = -1$ wird, und nach Bd. I, \S.\,171
\[
(r_1, \eta)^m = (-1, \eta)^m = \sqrt{p}^{\,m}
\]
ist $^1$).

Fassen wir das Ergebniss der S\"atze 3. bis 6. in eine Formel
zusammen, so ergiebt sich, wenn $\varepsilon$ ein
Einheitsfunctional, $\varphi$ irgend ein Functional des K\"orpers
$\Omega_m$ bedeutet, und $p$ eine Reihe von Primzahlen durchl\"auft,
die congruent mit 1 nach dem Modul $q$ sind, worunter dieselbe
Primzahl $p$ auch mehrmals vorkommen kann,
\weq{12}{\omega = (r, x_0)^m = \varepsilon \varphi^m \prod^{p} (r^{m_2}, \eta)^m.}

\noindent
$^1$) In des Verfassers Abhandlung in Bd. 8 der Acta mathematica ist
es vers\"aumt, den Fall des Satzes 5. besonders hervorzuheben.
